\begin{proof}
\textbf{1. Trivial (base) cases.}
\begin{itemize}
\item \textbf{Case $n=0$.}  
\[
(1+x)^{0}=1=1+0\cdot x .
\]

\item \textbf{Case $x=0$.}  
\[
(1+0)^{n}=1=1+n\cdot0 .
\]

\item \textbf{Case $x=-1$.}  Here $1+x=0$.
\begin{itemize}
\item If $n=0$, then $(1+x)^{0}=1=1-0$.
\item If $n=1$, then $0=1-1$.
\item If $n\ge2$, then $(1+x)^{n}=0$ while $1+nx=1-n\le0$, so $0\ge1-n$.
\end{itemize}
\end{itemize}

\medskip\noindent\textbf{2. Induction hypothesis.}  
Assume that for a fixed integer $n\ge1$,
\begin{align}
(1+x)^{n}\ge 1+nx \qquad\text{for all }x\ge-1 .
\tag{IH}
\end{align}

\medskip\noindent\textbf{3. Inductive step.}  
Let $x\ge-1$, whence $1+x\ge0$. Multiplying \eqref{IH} by the non‑negative factor $1+x$ preserves the inequality:
\begin{align}
(1+x)^{n}(1+x) &\ge (1+nx)(1+x) .
\tag{1}
\end{align}
The left‑hand side is $(1+x)^{n+1}$. Expanding the right‑hand side gives
\begin{align}
(1+nx)(1+x)=1+(n+1)x+nx^{2}.
\tag{2}
\end{align}
Substituting \eqref{2} into \eqref{1} we obtain
\begin{align}
(1+x)^{n+1}\ge 1+(n+1)x+nx^{2}.
\tag{3}
\end{align}
Since $n\ge1$ and $x^{2}\ge0$, the term $nx^{2}$ is non‑negative; dropping it can only decrease the right‑hand side:
\begin{align}
(1+x)^{n+1}\ge 1+(n+1)x .
\tag{4}
\end{align}
Equation \eqref{4} is Bernoulli’s inequality for the exponent $n+1$.

\medskip\noindent\textbf{4. Completion of the induction.}  
The inequality holds for the base case $n=0$, and the implication
\[
(1+x)^{n}\ge1+nx\;\Longrightarrow\;(1+x)^{n+1}\ge1+(n+1)x
\]
has been proved. By the principle of mathematical induction,  
\[
(1+x)^{n}\ge1+nx\qquad\text{for every integer }n\ge0\text{ and every }x\ge-1 .
\]

\medskip\noindent\textbf{5. Equality cases.}  
Equality may occur only at the steps where we did not lose a non‑negative term.

\begin{itemize}
\item In the base case $n=0$ we have $(1+x)^{0}=1=1+0x$ for all $x\ge-1$.
\item In the base case $x=0$ we have $(1+0)^{n}=1=1+n\cdot0$ for all $n\ge0$.
\item In the inductive step the discarded term is $nx^{2}$.  
If $n\ge1$ and $x\neq0$, then $nx^{2}>0$, so the inequality in \eqref{3} is strict and hence \eqref{4} is strict.  
If $x=0$, then $nx^{2}=0$ and equality is preserved for every $n$.  
If $x=-1$, then $1+x=0$ and $(1+x)^{n}=0$, $1+nx=1-n$. Equality holds exactly when $n=1$ (both sides equal $0$); for $n\ge2$ we have $0>1-n$.
\end{itemize}

Collecting these observations, equality holds precisely in the following situations:

\begin{tabular}{ll}
\toprule
Condition & Reason \\ \midrule
$n=0$ (any $x\ge-1$) & Both sides equal $1$.\\
$x=0$ (any $n\ge0$) & Both sides equal $1$.\\
$(x,n)=(-1,1)$ & Both sides equal $0$.\\ \bottomrule
\end{tabular}

No other pair $(x,n)$ with $x\ge-1$, $n\ge0$ yields equality.

\medskip\noindent\textbf{6. Conclusion.}  
For every real $x\ge-1$ and every integer $n\ge0$,
\[
(1+x)^{\,n}\ge 1+nx ,
\]
with equality exactly in the three cases listed above. \(\square\)
\end{proof}